% ---- score2 standard preamble (from: python3 score2_tools.py preamble) ----
\documentclass[10pt]{extarticle}
\usepackage[a4paper,margin=1.5cm]{geometry}
\usepackage[T1]{fontenc}
\usepackage{amsmath,amssymb}
\usepackage{graphicx}
\usepackage[export]{adjustbox}
\usepackage{multicol}
\usepackage{enumitem}
\usepackage{booktabs,array,tabularx}
\setlength{\columnsep}{0.7cm}
\setlength{\parindent}{0pt}
\setlength{\parskip}{3pt}
\pagestyle{empty}
\makeatletter
% figure registry (figdefs.tex is generated by the extract step)
\newcommand{\qfigdef}[4]{\@namedef{qf@file@#1}{#2}\@namedef{qf@w@#1}{#3}\@namedef{qf@wide@#1}{#4}}
\newcommand{\qf@check}[1]{\@ifundefined{qf@file@#1}{\PackageError{score2}{Unknown figure id #1}{}}{}}
\newcommand{\qf@img}[4][]{\includegraphics[width=#2\dimexpr\@nameuse{qf@w@#4}pt\relax,max width=#3,max height=0.45\textheight,#1]{\@nameuse{qf@file@#4}}}
% Full-width questions: finalize wraps every question that contains a wide figure in \qfullbegin...\qfullend,
% so the WHOLE question (stem, figures, options) is set across the page in reading order, then the two columns
% resume. The question's head - its number, stem and first figure (\qfullsplit marks the end of the head) -
% is kept on one page, so a number is never left at a page foot with its figure overleaf; the rest may break
% across pages normally, so no big gaps are left. Without a \qfullsplit, the whole question is kept together
% if it is under 45% of a page. Never write these yourself.
\newif\ifqf@full
\newif\ifqf@boxed
\newsavebox{\qf@box}
\newcommand{\qf@open}{\global\qf@fulltrue\par\global\qf@boxedtrue
  \setbox\qf@box\vbox\bgroup\hsize\textwidth\linewidth\textwidth}
\newcommand{\qf@emit}[1]{\par\egroup\global\qf@boxedfalse
  \ifdim\dimexpr\ht\qf@box+\dp\qf@box\relax<#1\textheight\noindent\usebox{\qf@box}\par\else\unvbox\qf@box\fi}
\newcommand{\qfullsplit}{\ifqf@boxed\qf@emit{0.85}\fi}
\newcommand{\qf@close}{\ifqf@boxed\qf@emit{0.45}\fi\global\qf@fullfalse}
\newcommand{\qfullbegin}{\end{multicols}\qf@open}
\newcommand{\qfullend}{\qf@close\begin{multicols}{2}}
\newcommand{\qfullstart}{\qf@open}
\newcommand{\qfulldone}{\qf@close}
% \qfig[scale]{id}: block figure, centred, directly where it is written (size set by the tool).
% Fallback if a wide figure is not inside a full-width question: break out of the columns in place - no floats.
\newcommand{\qfig}[2][1]{\qf@check{#2}%
  \ifqf@full\qf@blk{#1}{#2}\else
  \ifnum\@nameuse{qf@wide@#2}=1 \expandafter\@firstoftwo\else\expandafter\@secondoftwo\fi
  {\end{multicols}\begin{center}\qf@img{#1}{\textwidth}{#2}\end{center}\begin{multicols}{2}}%
  {\qf@blk{#1}{#2}}\fi}
\newcommand{\qf@blk}[2]{\par\smallskip{\centering\qf@img{#1}{\linewidth}{#2}\par}\smallskip}
% \qfigi[scale]{id}: inline figure (inside option lists / tables / minipages); never breaks columns.
\newcommand{\qfigi}[2][1]{\qf@check{#2}\qf@img[valign=c]{#1}{\linewidth}{#2}}
\makeatother
% structure macros - lint reads these to check completeness and figure placement
\newcommand{\papersec}[1]{\par\bigskip{\centering\Large\bfseries PAPER-#1\par}\medskip}
\newcommand{\partsec}[2]{\par\medskip\noindent{\large\bfseries PART-#1: #2}\par\smallskip}
\newcommand{\sect}[1]{\par\smallskip\noindent\textbf{#1}\par}
\newcommand{\qn}[1]{\par\smallskip\noindent\textbf{#1.}~}
\newcommand{\paraq}[2]{\par\smallskip\noindent\textbf{Paragraph for Questions #1 and #2}\par}
\newenvironment{qoptions}{\begin{itemize}[leftmargin=1.8em,itemsep=1pt,topsep=2pt,label={}]}{\end{itemize}}
\newcommand{\opt}[2]{\item[(#1)] #2}
% figure registry (generated by score2_tools.py finalize)
\qfigdef{p1_phy_s1_q01_1}{Pre-Score (2023-24) - Test 05 (05 Nov 2023)-img/p1_phy_s1_q01_1.png}{66.3}{0}
\qfigdef{p1_phy_s1_q03_1}{Pre-Score (2023-24) - Test 05 (05 Nov 2023)-img/p1_phy_s1_q03_1.png}{169.2}{0}
\qfigdef{p1_phy_s1_q04_1}{Pre-Score (2023-24) - Test 05 (05 Nov 2023)-img/p1_phy_s1_q04_1.png}{108.2}{0}
\qfigdef{p1_phy_s2_q05_1}{Pre-Score (2023-24) - Test 05 (05 Nov 2023)-img/p1_phy_s2_q05_1.png}{108.2}{0}
\qfigdef{p1_phy_s2_q06_1}{Pre-Score (2023-24) - Test 05 (05 Nov 2023)-img/p1_phy_s2_q06_1.png}{192.7}{0}
\qfigdef{p1_phy_s2_q07_1}{Pre-Score (2023-24) - Test 05 (05 Nov 2023)-img/p1_phy_s2_q07_1.png}{66.3}{0}
\qfigdef{p1_phy_s2_q08_1}{Pre-Score (2023-24) - Test 05 (05 Nov 2023)-img/p1_phy_s2_q08_1.png}{104.9}{0}
\qfigdef{p1_phy_s2_q09_1}{Pre-Score (2023-24) - Test 05 (05 Nov 2023)-img/p1_phy_s2_q09_1.png}{110.2}{0}
\qfigdef{p1_phy_s2_q10_1}{Pre-Score (2023-24) - Test 05 (05 Nov 2023)-img/p1_phy_s2_q10_1.png}{175.9}{0}
\qfigdef{p1_phy_s3_para1_2_1}{Pre-Score (2023-24) - Test 05 (05 Nov 2023)-img/p1_phy_s3_para1_2_1.png}{145.4}{0}
\qfigdef{p1_phy_s3_para3_4_1}{Pre-Score (2023-24) - Test 05 (05 Nov 2023)-img/p1_phy_s3_para3_4_1.png}{105.8}{0}
\qfigdef{p1_phy_s4_q07_1}{Pre-Score (2023-24) - Test 05 (05 Nov 2023)-img/p1_phy_s4_q07_1.png}{192.0}{0}
\qfigdef{p1_phy_s4_q08_1}{Pre-Score (2023-24) - Test 05 (05 Nov 2023)-img/p1_phy_s4_q08_1.png}{110.4}{0}
\qfigdef{p1_phy_s4_q09_1}{Pre-Score (2023-24) - Test 05 (05 Nov 2023)-img/p1_phy_s4_q09_1.png}{123.5}{0}
\qfigdef{p1_chem_s1_q01_A}{Pre-Score (2023-24) - Test 05 (05 Nov 2023)-img/p1_chem_s1_q01_A.png}{65.6}{0}
\qfigdef{p2_phy_s1_q01_1}{Pre-Score (2023-24) - Test 05 (05 Nov 2023)-img/p2_phy_s1_q01_1.png}{116.3}{0}
\qfigdef{p2_phy_s1_q02_1}{Pre-Score (2023-24) - Test 05 (05 Nov 2023)-img/p2_phy_s1_q02_1.png}{234.1}{0}
\qfigdef{p2_phy_s1_q03_1}{Pre-Score (2023-24) - Test 05 (05 Nov 2023)-img/p2_phy_s1_q03_1.png}{200.6}{0}
\qfigdef{p2_phy_s1_q04_1}{Pre-Score (2023-24) - Test 05 (05 Nov 2023)-img/p2_phy_s1_q04_1.png}{223.1}{0}
\qfigdef{p2_phy_s1_q05_1}{Pre-Score (2023-24) - Test 05 (05 Nov 2023)-img/p2_phy_s1_q05_1.png}{124.1}{0}
\qfigdef{p2_phy_s1_q05_2}{Pre-Score (2023-24) - Test 05 (05 Nov 2023)-img/p2_phy_s1_q05_2.png}{118.8}{0}
\qfigdef{p2_phy_s1_q06_1}{Pre-Score (2023-24) - Test 05 (05 Nov 2023)-img/p2_phy_s1_q06_1.png}{81.8}{0}
\qfigdef{p2_phy_s2_para7_8_1}{Pre-Score (2023-24) - Test 05 (05 Nov 2023)-img/p2_phy_s2_para7_8_1.png}{60.5}{0}
\qfigdef{p2_phy_s2_para9_10_1}{Pre-Score (2023-24) - Test 05 (05 Nov 2023)-img/p2_phy_s2_para9_10_1.png}{159.0}{0}
\qfigdef{p2_phy_s3_para3_4_1}{Pre-Score (2023-24) - Test 05 (05 Nov 2023)-img/p2_phy_s3_para3_4_1.png}{117.3}{0}
\qfigdef{p2_phy_s3_para5_6_1}{Pre-Score (2023-24) - Test 05 (05 Nov 2023)-img/p2_phy_s3_para5_6_1.png}{105.8}{0}
\qfigdef{p2_phy_s4_q07_1}{Pre-Score (2023-24) - Test 05 (05 Nov 2023)-img/p2_phy_s4_q07_1.png}{139.9}{0}
\qfigdef{p2_phy_s4_q08_1}{Pre-Score (2023-24) - Test 05 (05 Nov 2023)-img/p2_phy_s4_q08_1.png}{72.5}{0}
\qfigdef{p2_phy_s4_q09_1}{Pre-Score (2023-24) - Test 05 (05 Nov 2023)-img/p2_phy_s4_q09_1.png}{108.0}{0}
\qfigdef{p2_chem_s3_para3_4_1}{Pre-Score (2023-24) - Test 05 (05 Nov 2023)-img/p2_chem_s3_para3_4_1.png}{233.5}{0}
\qfigdef{p2_chem_s3_q03_1}{Pre-Score (2023-24) - Test 05 (05 Nov 2023)-img/p2_chem_s3_q03_1.png}{184.7}{0}
% ---- body skeleton ----
% \begin{document}
% \begin{center}{\Large\bfseries JEE (Main+Advanced) Score-II --- Test NN (DD Mon YYYY)}\end{center}
% \begin{multicols}{2}
%   \papersec{1}   (only for two-paper tests)
%   \partsec{1}{PHYSICS} \sect{Section-I (i): Single Correct} \qn{1} ... \qfig{p1_phy_q01_1} ...
% \end{multicols}
% \end{document}

\begin{document}
\begin{center}{\Large\bfseries Pre-Score (2023-24) --- Test 05 (05 Nov 2023)}\end{center}\begin{multicols}{2}
\papersec{1}
\partsec{1}{PHYSICS}
\sect{SECTION-I(i)}
\qn{1} If we find the moment of inertia of a scalene triangular plate ABC about its base as shown. The mass of the plate is $M$. The base length is $b$ and altitude is $h$. Moment of inertia is :-
\qfig{p1_phy_s1_q01_1}
\begin{qoptions}
\opt{A}{$\dfrac{Mh^2}{12}$}
\opt{B}{$\dfrac{Mh^2}{3}$}
\opt{C}{$M\dfrac{h^2+b^2}{12}$}
\opt{D}{$\dfrac{Mh^2}{6}$}
\end{qoptions}
\qn{2} Assume that earth has a mass of $6\times10^{24}$ kg and a radius of 6400 km. Earth has a population of 8 billion people each of mass 50 kg. On the call of our prime minister, the entire population of earth lines up on the equator and move with a velocity of 1 m/s relative to the earth from west to east. What is the change in the time period of revolution of earth?
\begin{qoptions}
\opt{A}{Approx. $10^{-15}$ sec}
\opt{B}{Approx. $10^{-10}$ sec}
\opt{C}{Approx. $10^{-12}$ sec}
\opt{D}{Approx. $10^{-17}$ sec}
\end{qoptions}
\qn{3} In figure 1, we see a ball C colliding against a ball B which is connected to a ball A by a loose inextensible string. There is no friction anywhere. In figure 2, the ball C collides against ball B which is connected to ball A by a taut inextensible string. If the collision is elastic in each of the cases,
\qfig{p1_phy_s1_q03_1}
\begin{qoptions}
\opt{A}{The mechanical energy is conserved during and after the collision in both the cases.}
\opt{B}{The mechanical energy is conserved during and after the collision in case 1 but not in 2.}
\opt{C}{The mechanical energy is conserved during and after the collision in case 2 but not in 1.}
\opt{D}{The mechanical energy is not conserved during and after the collision in both the cases.}
\end{qoptions}
\qn{4} An experienced pool player can control the placement of a ball after it hits an edge bumper by putting spin, or ``English,'' on the ball. This is accomplished by aiming the cue tick so that it hits the ball either left or right of the ball's centre. Suppose a player gives a ball English by hitting it to the left of centre (left from his perspective). Which path in Figure does the ball take after it collides with the bumper? How is it rotating before collision with the bumper as seen from the top?
\qfig{p1_phy_s1_q04_1}
\begin{qoptions}
\opt{A}{Path a with clockwise spin.}
\opt{B}{Path b with clockwise spin.}
\opt{C}{Path c with anticlockwise spin.}
\opt{D}{Path a with anticlockwise spin.}
\end{qoptions}
\sect{SECTION-I(ii)}
\qn{5} A uniform circular ring of mass $m=1$kg and radius $R=20$ cm is placed on a rough inclined plane of inclination $\theta=37^\circ$ as shown in the figure. Then choose the correct option(s) just after the ring is released from rest. (there is no slipping between the ring and inclined plane and $g=10$ m/s$^2$)
\qfig{p1_phy_s2_q05_1}
\begin{qoptions}
\opt{A}{The frictional force acting on the ring due to inclined plane is 4N.}
\opt{B}{The frictional force acting on the ring due to inclined plane is 3N.}
\opt{C}{The normal force acting on the ring due to inclined plane is 8 N.}
\opt{D}{The normal force acting on the ring due to inclined plane is 11N.}
\end{qoptions}
\qn{6} The 40-kg package is thrown with a speed of 4 m/s onto the cart having a mass of 20 kg. If it slides on the smooth surface and strikes the spring, determine the velocity $u_c$ of the cart at the instant the package fully compresses the spring. What is the maximum compression $x$ of the spring? Neglect rolling resistance of the cart. What is the velocity $v_p$ of the package after the spring comes back to the natural length? What is the velocity $v_c$ of the cart at that instant?
\qfig{p1_phy_s2_q06_1}
\begin{qoptions}
\opt{A}{$u_c=\dfrac{4}{\sqrt3}$ m/s}
\opt{B}{$v_c=\dfrac{8}{\sqrt3}$ m/s}
\opt{C}{$v_p=\dfrac{2}{\sqrt3}$ m/s}
\opt{D}{$x=\sqrt{\dfrac{2}{75}}$ m}
\end{qoptions}
\qn{7} A small uniform ring of mass $m$ and radius $R$ is held at the corner of a table and released. It was in unstable equilibrium and so it falls off to the ground below. We have two options for the table, both having the same height. The first table has a rough edge so that the ring does not slip when it is in contact with the edge, and the second table has a smooth edge.
\qfig{p1_phy_s2_q07_1}
\begin{qoptions}
\opt{A}{In the case of the first table, the ring rotates by a larger angle before leaving contact with the edge.}
\opt{B}{In the case of first table, the angular velocity of the ring remains constant during its motion in the air.}
\opt{C}{In the case of second table, the ring travels with lesser horizontal velocity during its motion in the air.}
\opt{D}{In the case of second table, angular velocity of ring in air is more than ring first table.}
\end{qoptions}
\qn{8} If the system is released from rest, determine the speed of the 20-kg cylinders A and B after A has moved by a distance of 2 m. The differential pulley has a mass of 15 kg with a radius of gyration about its centre of mass $k=100$ mm.
\qfig{p1_phy_s2_q08_1}
\begin{qoptions}
\opt{A}{A will move up.}
\opt{B}{B will move up.}
\opt{C}{A will have a speed of $\dfrac{4\sqrt{15}}{5}$ m/s.}
\opt{D}{B will have a speed of $\dfrac{30}{\sqrt{2445}}$ m/s.}
\end{qoptions}
\qn{9} A uniform cylinder of mass $M$, height $h$ and diameter $d$ is kept on a rough inclined plane. The coefficient of static friction and kinetic friction are both $\mu$. Select the correct statement(s).
\qfig{p1_phy_s2_q09_1}
\begin{qoptions}
\opt{A}{If $\mu>\tan\beta$ and $\tan\beta>d/h$ Cylinder will neither topple nor slide.}
\opt{B}{If $\mu>\tan\beta$ and $\tan\beta>d/h$ Cylinder will topple but will not slide.}
\opt{C}{If $\mu<\tan\beta$ and $\tan\beta>d/h$ Cylinder will slide but not topple.}
\opt{D}{If $\mu<\tan\beta$ and $\tan\beta<d/h$ Cylinder will slide but not topple.}
\end{qoptions}
\qn{10} At the instant shown both rods of negligible mass swing with a counterclockwise angular velocity of 5 rad/s while the 50 kg bar is subjected to the 100 N horizontal force. Determine the tension developed in the rods and the angular acceleration of the rods at this instant.
\qfig{p1_phy_s2_q10_1}
\begin{qoptions}
\opt{A}{The angular acceleration is anti-clockwise.}
\opt{B}{The angular acceleration is $4/3$ rad/s$^2$.}
\opt{C}{Tension is approx. 1.18 kN.}
\opt{D}{Tension is more in rod AB than in CD.}
\end{qoptions}
\sect{SECTION-II(i)}
\paraq{1}{2}
A wheel of radius $R=10.0$ cm moves without slippage along horizontal surface with a constant velocity $v_0=2.0$ m/s.
\qfig{p1_phy_s3_para1_2_1}
\qn{1} Determine maximum height $H$ (in cm) that can be reached by a splash of mud, which is separated from that moving wheel. Assume that friction with air can be neglected, while height $H$ is measured from the horizontal surface. For calculations use acceleration due to gravity as 10 m/s$^2$.
\qn{2} Find the maximum height (in cm) from the ground reached by the splash of mud that just grazes the moving wheel again when it is travelling downwards.
\paraq{3}{4}
The 20-kg block A is towed up the ramp of the 40-kg cart using the motor M mounted on the side of the cart. The motor begins winding in the cable with a constant velocity of 5m/s, measured relative to the cart. Both the block and cart are at rest initially. The kinetic friction between the block and the ramp is absent. Neglect rolling resistance.
\qfig{p1_phy_s3_para3_4_1}
\qn{3} Determine how far (in m) the cart will move when the block has travelled a distance 2 m up the ramp.
\qn{4} Determine the work done by the motor (in J) till then.
\paraq{5}{6}
A uniform rod of mass 100 g and length 50.0 cm rotates in a horizontal plane about a fixed, vertical, frictionless pin through its centre. Two small beads, each of mass 30.0 g, are mounted on the rod so that they are able to slide without friction along its length. Initially the beads are held by catches at positions 10.0 cm on each side of centre, at which time the system rotates at an angular speed of 20.0 rad/s. Suddenly, the catches are released and the small beads slide outward along the rod.
\qn{5} Find the angular speed of the system (in rad/s) at the instant the beads reach the ends of the rod.
\qn{6} If the beads fly off the ends, What is the speed (in m/s) of each of the beads just after that? Fill $v^2$.
\sect{SECTION-II(ii)}
\qn{7} The sports car has a mass of 1500 kg and a centre of mass at G. Determine the shortest time (in sec) it takes for it to reach a speed of 108 km/hr, starting from rest, if the engine only drives the rear wheels, whereas the front wheels are free rolling. The coefficient of static friction between the wheels and the road is 0.2. Neglect the mass of the wheels for the calculation. Round off to nearest integer.
\qfig{p1_phy_s4_q07_1}
\qn{8} The disk has a mass of 20 kg and is originally spinning at the end of the massless rod with an angular velocity of $1000(2-\sqrt3)$ rad/s. If it is then placed against the wall, where the coefficient of kinetic friction is $2-\sqrt3$, determine the time (in sec) required for the motion to stop.
\qfig{p1_phy_s4_q08_1}
\qn{9} A ring with uniformly distributed mass $m$ hangs at the fixed horizontal cylinder, which has radius slightly less than inner radius of the ring $R$. Coefficient of friction between the ring and fixed cylinder is $\mu=0.1$. On top of the ring is thrown a non-stretchable rope with negligible mass, where weights with masses $m_A=2m$ and $m_B=5m$ are attached to the ends of the rope.
Determine acceleration (in m/s$^2$) $a$ of the weight B after release of the system in the gravity of the Earth. For calculations use acceleration due to gravity as $g=10$ m/s$^2$. Assume that rope does not slip on the ring. Round off to nearest integer.
\qfig{p1_phy_s4_q09_1}
\partsec{2}{CHEMISTRY}
\sect{SECTION-I(i)}
\qn{1} Which of the following species cannot be resolved into `d' and `$\ell$' form.
\begin{qoptions}
\opt{A}{\qfigi{p1_chem_s1_q01_A}}
\opt{B}{[Pt(bn)$_2$]Cl$_2$}
\opt{C}{Na[Co(CN)$_3$(NH$_3$)$_2$Cl]}
\opt{D}{[Zn(gly)$_2$]}
\end{qoptions}
\qn{2} Sodium lauryl sulphate is widely used as a collector in froth flotation process. It can be written as C$_{12}$H$_{25}$OSO$_3$Na. Which is correct regarding its function as a collector particle during froth floatation.
\begin{qoptions}
\opt{A}{Its hydrophillic end gets attached to sulphide ore particle which then gets attached to air bubble}
\opt{B}{Its hydrophobic end tends to stay in water thereby facilitating ore particle movement towards surface}
\opt{C}{It stays at air-water interface and the ore settles at the bottom}
\opt{D}{It's hydrophobic end gets attached to the ore particle and thus enhances the floating characteristics}
\end{qoptions}
\qn{3} A colourless odourless gas ``X'' can be prepared by action of steam over hot coke on a surface of catalyst at high temperature. The same gas ``X: reduces 2-ethylanthraquinone in the presence of Ni/Pd catalyst. Which of the following is \textbf{CORRECT} for gas ``X''.
\begin{qoptions}
\opt{A}{Bond strength of F$_2$ $>$ Bond strength of X}
\opt{B}{``X'' is obtained at cathode in DOWN's cell}
\opt{C}{``X'' boils at lower temperature as compared to lightest noble gas}
\opt{D}{``X'' react with iodine only in the presence of catalyst.}
\end{qoptions}
\qn{4} Which of the following metals involve oxidation of metal followed by reduction during their extraction.
\begin{qoptions}
\opt{A}{Ag from Ag$_2$S}
\opt{B}{Fe from Fe$_2$O$_3$}
\opt{C}{Au from Auriferrous rocks}
\opt{D}{Sn from SnO$_2$}
\end{qoptions}
\sect{SECTION-I(ii)}
\qn{5} Metal ``M'' can be reduced from its oxide via carbon reduction. It produces H$_2$ on dissolving metal in aqueous alkali. The complex MCl$_2$.4NH$_3$ produces white precipitate with AgNO$_3$. If M and M$^{2+}$ are diamagnetic in gaseous state, then
\begin{qoptions}
\opt{A}{``M'' is refined by distillation}
\opt{B}{``M'' is refined by liquation}
\opt{C}{Electrorefining of ``M'' is done using aqueous electrolyte}
\opt{D}{One of the ore of M can be concentrated by froth floatation}
\end{qoptions}
\qn{6} Which of the following complex specie(s) are paramagnetic \& low spin.
\begin{qoptions}
\opt{A}{[NiF$_6$]$^{2-}$}
\opt{B}{[Fe(CN)$_6$]$^{3-}$}
\opt{C}{[Co(CN)$_6$]$^{4-}$}
\opt{D}{[Cr(CN)$_6$]$^{3-}$}
\end{qoptions}
\qn{7} Increase in concentration of which specie(s) will cause breakdown of ozone in atmosphere.
\begin{qoptions}
\opt{A}{Cl$^{\bullet}$}
\opt{B}{ClO$^{\bullet}$}
\opt{C}{O$_2$}
\opt{D}{NO$_2$ClO}
\end{qoptions}
\qn{8} ``X'' is used as solvent for dry cleaning and is carcinogenic in nature. ``Y'' is used as a bleaching agent in process of laundry. ``X'' is planar ``Y'' is non-planar.\par
Which is/are correct for ``X'' \& ``Y''.
\begin{qoptions}
\opt{A}{Both ``X'' \& ``Y'' are non-polar}
\opt{B}{``X'' is non polar \& ``Y'' is polar}
\opt{C}{Urea acts as a stabilizer to store ``Y''}
\opt{D}{Liquified CO$_2$ can be used for dry-cleaning in place of ``X''}
\end{qoptions}
\qn{9} Which is/are \textbf{CORRECT}.
\begin{qoptions}
\opt{A}{CuSO$_4$.5H$_2$O contains hydrogen bonded H$_2$O.}
\opt{B}{BaCl$_2$.2H$_2$O is interstitial water}
\opt{C}{FeSO$_4$.7H$_2$O has six H$_2$O are coordinated with metal}
\opt{D}{Fluorine reacts with NaOH(aq.) to produce OF$_2$}
\end{qoptions}
\qn{10} In given four options, each option contains a complex salt which have cationic as well as anionic part is complex sphere. Which of the following coordination compound show co-ordination isomerism as well as stereo isomerism.
\begin{qoptions}
\opt{A}{[Pt(en)$_2$] [PtCl$_4$]}
\opt{B}{[Co(NH$_3$)$_6$] [Co(CN)$_4$(NO$_2$)$_2$]}
\opt{C}{[Pt(en)$_2$Br$_2$] [PtCl$_4$]}
\opt{D}{[Ag(NH$_3$)$_2$] [Ag(CN)$_2$]}
\end{qoptions}
\sect{SECTION-II(i)}
\paraq{1}{2}
Effective atomic number provides a convenient way to compare stability of carbonyl complexes. It is related to Sidgwick's EAN rule which compared complex compounds using analogy of inert gases.
\qn{1} The sum of magnetic moments of the complexes that do not follow Sidwick's EAN rule among [Fe(CO)$_5$] , K$_2$[PtCl$_6$] , [Fe(NH$_3$)$_6$]Cl$_2$ , [Ni(NH$_3$)$_6$]SO$_4$ is \underline{\hspace{1cm}}.
\qn{2} Number of complex compounds from given below which have more number of valence electrons in the orbital of each central metal as compared to the number of valence electrons on central metal in [Mn(CO)$_5$] complex.\par
Mn$_2$(CO)$_{10}$ , \ [Fe(CO)$_5$] , \ [Ni(CO)$_4$] , \ Co$_2$(CO)$_8$ , \ Fe$_2$(CO)$_9$
\paraq{3}{4}
Leaching is a chemical method for concentration of ore involving a leaching agent and other reagent. Usually ores of Au, Al , Ag, Cu are concentrated by this method.
\qn{3} During leaching of Bauxite by Serpeck's method, two gaseous compounds are formed. The sum of oxidation states of all the least electronegative atom(s) in both compounds is \underline{\hspace{1cm}}.
\qn{4} Leaching of Ag in dilute KCN in presence of O$_2$ produces a compound X. Sum of coordination number of X and EAN of X is \underline{\hspace{1cm}}.
\paraq{5}{6}
Tropospheric pollution occurs due to presence of undesirable solids or gaseous particles in air. These pollutants are divided into two categories viz. gaseous pollutants and particulate pollutants.
\qn{5} Number of reaction(s) that yield a gas which can act as gaseous air pollutant.\par
(i) NH$_3$ + O$_2$ $\xrightarrow[\text{500K, 9 bar}]{\text{Pt--Rh catalyst}}$\par
(ii) NaSO$_3$ $\xrightarrow{\text{Zn + dil. H}_2\text{SO}_4}$\par
(iii) NaNO$_2$ $\xrightarrow{\text{dil. HCl}}$\par
(iv) HCOOH $\xrightarrow{\text{conc. H}_2\text{SO}_4}$
\qn{6} To create a brown haze in an artificial atmosphere, 5 mole NO$_2$ is radiated with energy equivalent to that of sunlight to break NO$_2$ molecule into NO \& [O] atom.\par
Sufficient O$_2$ is supplied for formation of O$_3$ by reaction of O$_2$ and [O]. If percentage yield of dissociation of NO$_2$ is 80\% , then calculate amount of NO$_2$ finally present in artificial atmosphere (in moles)
\sect{SECTION-II(ii)}
\qn{7} A sample of hard water which contain equimolar mixture of Ca(HCO$_3$)$_2$ and Mg(HCO$_3$)$_2$ , is treated with 3000 mL, 0.5 M Ca(OH)$_2$ solution such that none of the reactants are left. Amount of precipitate formed (in grams) will be \underline{\hspace{1cm}}.\par
(Consider all Mg$^{2+}$ to be precipitated as Mg(OH)$_2$ and all Ca$^{2+}$ as CaCO$_3$)\par
(Ca = 40, C = 12, Mg = 24, O = 16, H = 1)
\qn{8} Number of linkage isomers of Na$_4$[Fe(CN)$_5$(NOS)] will be :-
\qn{9} Number of process(es) that involve no chemical change\par
(i) Distillation of Zn\par
(ii) Mond's process\par
(iii) Liquation\par
(iv) Poling\par
(v) Zone refining\par
(vi) Thermal decomposition\par
(vii) Mac-Arthur forest process\par
(viii) Cupellation\par
(ix) Amalgamation\par
(x) Parke's method
\partsec{3}{MATHEMATICS}
\sect{SECTION-I(i)}
\qn{1} Let $M_n$ be the $n\times n$ matrix with entries as follow for $1\le i\le n$, $m_{i,i}=10$; for $1\le i\le n-1$, $m_{i+1,i}=m_{i,i+1}=3$ all other entries in $M_n$ are zero. Let $D_n$ be the determinant of matrix $M_n$, then $\displaystyle\sum_{n=1}^{\infty}\dfrac{1}{8D_n+1}$
\begin{qoptions}
\opt{A}{$\dfrac14$}
\opt{B}{$\dfrac18$}
\opt{C}{$\dfrac1{36}$}
\opt{D}{$\dfrac1{72}$}
\end{qoptions}
\qn{2} Let $\omega$ be a complex cube root of unity with $\omega\ne1$ and $P=[p_{ij}]$ be a $n\times n$ matrix with $p_{ij}=\omega^{i+j}$. Then $P^2=0$, when $n=$
\begin{qoptions}
\opt{A}{57}
\opt{B}{55}
\opt{C}{58}
\opt{D}{56}
\end{qoptions}
\qn{3} Evaluate $\displaystyle\sum_{k=0}^{n+1}\sum_{j=1}^{k+1}\sum_{i=1}^{k}{}^iC_{j-1}$ where $n>2$ and ${}^nC_r=0\ \forall\ r>n$
\begin{qoptions}
\opt{A}{$2^{n+3}-2n-6$}
\opt{B}{$2^{n+3}-2n-4$}
\opt{C}{$2^{n+2}-4n-2$}
\opt{D}{None of these}
\end{qoptions}
\qn{4} Let $x$ be a fixed real number which is not an integer multiple of $\pi$. For each natural number $n$ let $A_n=[a_{jk}]$ be the $n\times n$-matrix defined by\par
\[a_{jk}=\begin{cases}0 & ;\ \text{for } |j-k|>1\\ 1 & ;\ \text{for } |j-k|=1\\ 2\cos x & ;\ \text{for } j=k\end{cases}\]
If $D_n=\det A_n$ then the solution of $D_3-2D_2\cos x+D_1=0$ is $[0,2\pi]$ is
\begin{qoptions}
\opt{A}{infinitely many}
\opt{B}{0}
\opt{C}{2}
\opt{D}{3}
\end{qoptions}
\sect{SECTION-I(ii)}
\qn{5} The number of ways to paint a regular octahedron (it has eight equilateral triangular faces, 6 vertices, 12 sides) with 8 different colors is (All colors to be used)
\begin{qoptions}
\opt{A}{Same as the indistinguishable number of ways to name the vertices of a cube with letters A, B, C, D, E, F, G, H}
\opt{B}{having sum of digits 15}
\opt{C}{having 32 positive divisors}
\opt{D}{having four prime factors}
\end{qoptions}
\qn{6} If the elements of a $2\times2$ matrix $A$ are real positive and distinct such that $\det(A+A^T)=0$ where $\det(A)$ represent determinant value of matrix $A$, then :
\begin{qoptions}
\opt{A}{$\det A>0$}
\opt{B}{$\det A\ge0$}
\opt{C}{$\det(A-A^T)>0$}
\opt{D}{$\det(AA^T)>0$}
\end{qoptions}
\qn{7} If $A$ is non-singular matrix of order $n\times n$ such that $3ABA^{-1}+A=2A^{-1}BA$, then :
\begin{qoptions}
\opt{A}{$A$ and $B$ both are identity matrices}
\opt{B}{$|A+B|=0$}
\opt{C}{$|ABA^{-1}-A^{-1}BA|=0$}
\opt{D}{$A+B$ is not a singular matrix}
\end{qoptions}
\qn{8} Let `$n$' be an integer with $n\ge3$. Let $P_1P_2P_3\ldots P_n$ be a regular $n-$sided polygon inscribed in a circle. Three points $P_i$, $P_j$ and $P_k$ are randomly chosen, where $i$, $j$ and $k$ are distinct integers from 1 to $n$. The number of triangles such that the triangle $P_iP_jP_k$ is obtuse is
\begin{qoptions}
\opt{A}{$\dfrac{n(n-2)(n-4)}{8}$, if $n$ is even}
\opt{B}{$\dfrac{n(n-1)(n-3)}{8}$, if $n$ is odd}
\opt{C}{$\dfrac{n(n-1)(n-2)}{6}$, if $n$ is even}
\opt{D}{$\dfrac{n(n+1)(n+2)}{6}$, if $n$ is odd}
\end{qoptions}
\qn{9} Let $N$ be the number of ways to place the integers 1 through 12 in the 12 cells of a $2\times6$ grid so that for any two cells sharing a side, the difference between the numbers in those cells is not divisible by 3. One way to do this is shown below.
\begin{center}
\renewcommand{\arraystretch}{1.4}\setlength{\tabcolsep}{5pt}\footnotesize
\begin{tabular}{|c|c|c|c|c|c|}\hline
1&3&5&7&9&11\\\hline
2&4&6&8&10&12\\\hline
\end{tabular}
\end{center}
\begin{qoptions}
\opt{A}{$N=2^{11}\cdot3^5\cdot5$}
\opt{B}{$N=2^{11}\cdot3^5\cdot5^2$}
\opt{C}{Total divisor of $N=144$}
\opt{D}{Total divisor of $N=216$}
\end{qoptions}
\qn{10} If value of $3P={}^{30}C_5-{}^{30}C_6+{}^{30}C_7-\ldots-{}^{30}C_{28}+29$, and value of ${}^QC_R={}^{100}C_6+4\,{}^{100}C_7+6\,{}^{100}C_8+4\,{}^{100}C_9+{}^{100}C_{10}$, then
\begin{qoptions}
\opt{A}{number of positive divisors of $P$ are 16}
\opt{B}{$Q+R$ can be divisible either by 9 or 19}
\opt{C}{if $abcd=P$, $a,b,c,d\in I^+$, then possible number of order pairs of $(a,b,c,d)$ are $4^4$}
\opt{D}{$P$ is an even integer}
\end{qoptions}
\sect{SECTION-II(i)}
\paraq{1}{2}
Consider two $3\times3$ matrices $A$ and $B$ satisfying $A=\text{adj}(B)-B^T$ and $B=\text{adj}(A)-A^T$. Also given that $A$ is a non singular matrix. (Where $P^T$ denotes transpose of matrix $P$).
\qn{1} $|A|+|B|$ is equal to :
\qn{2} $|AB+BA|$ is equal to :
\paraq{3}{4}
If $(1+px+x^2)^n=1+a_1x+a_2x^2+\ldots+a_{2n}x^{2n}$ , ($n\in N$) then
\qn{3} The remainder obtained, when $a_1+5a_2+9a_3+7a_4+\ldots\ldots+(8n-3)a_{2n}$, is divided by $(p+2)$ is :
\qn{4} The value of $a_1+3a_2+5a_3+7a_4+\ldots\ldots+(4n-1)a_{2n}$, when $p=-3$ and $n=4$ is :
\paraq{5}{6}
A triangle is called an integral triangle if all its sides are integers. If $a$, $b$, and $c$ are sides of an integral triangle, then we an assume that $a\le b\le c$, because any other permutation of $a$, $b$, and $c$ will yield the same result. Now, answer the following question.
\qn{5} The number of integral isosceles or equilateral triangle's, none of whose side exceeds 4, is :
\qn{6} The number of integral isosceles or equilateral triangles, none of whose side exceeds 6 is :
\sect{SECTION-II(ii)}
\qn{7} Ram chooses five subsets $S_1,S_2,S_3,S_4,S_5$ of $\{1,2,3,4,5,6,7\}$ such that $|S_1|=1$, $|S_2|=2$, $|S_3|=3$, $|S_4|=4$, and $|S_5|=5$. Moreover, for all $1\le i<j\le5$, either $S_i\cap S_j=S_i$ or $S_i\cap S_j=\phi$ (in other words, the intersection of $S_i$ and $S_j$ is either $S_i$ or the empty set). If Ram can achieve this in $\lambda$ ways then the $3^{rd}$ digit of $\lambda$ from begining is \underline{\hspace{1.5cm}}
\qn{8} If $p=\displaystyle\sum_{r=1}^{50}\left\{\sum_{k=1}^{r}\left(\dfrac{(-1)^{r-1}}{k}\cdot{}^{50}C_r\right)\right\}$ then the value of $200(p)$ equals
\qn{9} Let $A=\begin{bmatrix}a&b\\c&d\end{bmatrix}$ where $a,b,c,d$ are positive integers arranged in ascending order and exactly two of $a,b,c,d$ are prime numbers and also pair wise coprime and satisfies $A^2-16A-17I=0$ (where $I$ is an identity matrix of order $2\times2$). If $B=\begin{bmatrix}b&a\\c&d\end{bmatrix}$ where $\det(A)$ represent determinant value of matrix $A$, then Det $B$ equals
\papersec{2}
\partsec{1}{PHYSICS}
\sect{SECTION-I(i)}
\qn{1} A semi-annular disc of mass $m$ has dimensions as shown. Find the moment of inertia of the disc about different axes.
\qfig{p2_phy_s1_q01_1}
\begin{qoptions}
\opt{A}{$I_x=\dfrac{M\left(R_1^2+R_2^2\right)}{4}$}
\opt{B}{$I_y=\dfrac{M\left(R_1^2+R_2^2\right)}{4}$}
\opt{C}{$I_z=\dfrac{M\left(R_1^2+R_2^2\right)}{2}$}
\opt{D}{$I$ about any axis in the $x$-$y$ plane passing through the origin $=\dfrac{M\left(R_1^2+R_2^2\right)}{4}$}
\end{qoptions}
\qn{2} A 4 kg projectile travels with a horizontal velocity of 600 m/s before it explodes and breaks into two fragments A and B of mass 1.5 kg and 2.5 kg, respectively. If the fragments travel along the parabolic trajectories shown,
\qfig{p2_phy_s1_q02_1}
\begin{qoptions}
\opt{A}{Both the fragments hit the ground simultaneously}
\opt{B}{A reaches ground after B}
\opt{C}{$V_A=800\,(\sqrt6+\sqrt2)$ m/s}
\opt{D}{$V_B=960\,(\sqrt3+1)$ m/s}
\end{qoptions}
\qn{3} A 100 gm rocket is fired vertically. The engine applies a thrust $\mathbf P$ which varies in magnitude as shown. Neglecting air resistance and the change in mass of the rocket, (Take : $g=10$ m/s$^2$)
\qfig{p2_phy_s1_q03_1}
\begin{qoptions}
\opt{A}{the maximum speed of the rocket as it goes up is 39 m/s.}
\opt{B}{the time for the rocket to reach its maximum elevation is 4.7 sec.}
\opt{C}{The total impulse of the engine is 4.7 Ns.}
\opt{D}{the impulse of gravity during the time engine acts is 0.8 Ns.}
\end{qoptions}
\qn{4} Figure 1 shows the block of mass $m$ connected to a disc of mass $M$ as shown. The disc can rotate freely about its massless axle and the string does not slip on the disc. Figure 2 shows the same disc hanging from an inextensible thread. The string does not slip on the disc. The systems are released from rest simultaneously. The tension in both the strings is $F$ as shown.
\qfig{p2_phy_s1_q04_1}
\begin{qoptions}
\opt{A}{$m=M$}
\opt{B}{In time $t$, kinetic energy of disc in both the cases is the same}
\opt{C}{In time $t$, the rotational kinetic energy of both the discs is same}
\opt{D}{The angular acceleration of disc is same in both the cases.}
\end{qoptions}
\qn{5} We have a fixed hemispherical surface of radius $R$. We place a sphere and a ring on the hemisphere one by one and release them. Both of them roll on the surface till they leave the surface.
\begin{center}\qfigi{p2_phy_s1_q05_1}\hfill\qfigi{p2_phy_s1_q05_2}\end{center}
\begin{qoptions}
\opt{A}{The Sphere covers a larger angle $\theta$ on the hemisphere.}
\opt{B}{The ring has higher kinetic energy at the instant of leaving the hemisphere.}
\opt{C}{The sphere has higher speed at the instant of leaving the hemisphere.}
\opt{D}{The work done by contact force is same in both the cases.}
\end{qoptions}
\qn{6} A ball is projected towards a smooth vertical wall 3m away. It collides the wall \& lands back at the point of projection. Choose the correct statement(s) :-
\qfig{p2_phy_s1_q06_1}
\begin{qoptions}
\opt{A}{time taken to reach wall is more than time taken to reach back after collision.}
\opt{B}{time taken to reach wall is less than time taken to reach back after collision.}
\opt{C}{the coefficient of restitution is $\dfrac5{11}$.}
\opt{D}{the coefficient of restitution is $\dfrac13$}
\end{qoptions}
\sect{SECTION-I(ii)}
\paraq{7}{8}
Children try to push the door at any point to open it. But the adult can easily open the door by pushing the door as far from the hinge as possible. That is why we see the handle of the door at the far end, away from the hinge. While the idea can be easily explained by concept of torque and the lever arm, a question arises that if the same force is applied at a point close to the hinge, acceleration of center of mass of door is lesser. Which force reduces this effect ? We will try to answer this question . Consider a uniform rectangular door of height $h$ and width $w$, mass $m$. Hinges are smooth. Neglect the effect of gravity.
\qfig{p2_phy_s2_para7_8_1}
\qn{7} A force $F$ is applied perpendicular to the door at a distance of $\dfrac w3$ from hinge. What is the force exerted by the hinge at this moment ?
\begin{qoptions}
\opt{A}{0.5 $F$ in direction opposite to the force.}
\opt{B}{0.5 $F$ in direction of the force.}
\opt{C}{$\dfrac F3$ in direction opposite to the force}
\opt{D}{$\dfrac{2F}3$ in direction opposite to the force.}
\end{qoptions}
\qn{8} How far from the axis should the force be applied so that force exerted by hinge is zero at that moment.
\begin{qoptions}
\opt{A}{$\dfrac w2$}
\opt{B}{$\dfrac{3w}4$}
\opt{C}{$\dfrac{2w}3$}
\opt{D}{$\dfrac{7w}{12}$}
\end{qoptions}
\paraq{9}{10}
The uniform plate has a weight of 500 N. Determine the tension in each of the supporting cables. There are only 3 cables supporting the rectangular plate. A force of 200 N is also acting at the mid-point of an edge as shown in figure.
\qfig{p2_phy_s2_para9_10_1}
\qn{9} Tension in cable C is
\begin{qoptions}
\opt{A}{100 N}
\opt{B}{200 N}
\opt{C}{400 N}
\opt{D}{500 N}
\end{qoptions}
\qn{10} Tension in cable B is
\begin{qoptions}
\opt{A}{same as that of A but less than that in C}
\opt{B}{more than that of A but less than that in C}
\opt{C}{less than that of A but more than C}
\opt{D}{same as that of A and more than that in C}
\end{qoptions}
\sect{SECTION-II(i)}
\paraq{1}{2}
Two particles of mass 1kg and 4 kg are moving in empty space with velocity of 1m/s in negative $x$ direction and 4 m/s in positive $x$ direction respectively. 1 kg mass is at the origin of coordinates at the given instant and 4 kg mass is at $x=4$m at that instant. They exert an attractive force on each other which depends on the distance $s$ between them as $F=5s$.
\qn{1} What is the maximum distance $s_{max}$ (in m) between them during the course of their motion. Fill $(2s_{max})^2$ ?
\qn{2} What is the distance $s$ (in m) between them when 1 kg mass stops for the first time. Fill $(2s)^2$
\paraq{3}{4}
Tennis balls are usually rejected if they fail to rebound to waist level when dropped from shoulder level. A ball just passes the test as indicated in the figure. Assume that the coefficient of restitution between the ball and the floor is the same as the coefficient of restitution between the balls.
\qfig{p2_phy_s3_para3_4_1}
\qn{3} What is the coefficient of restitution for the collision between the floor and the ball.
\qn{4} If two identical tennis balls moving at the same speed in opposite direction collide head on, what is the percentage of total mechanical energy that is lost?
\paraq{5}{6}
A jet of fluid with cross-sectional area $A$ and mass density $\rho$ issues from the nozzle with a velocity $v$ and impinges on the inclined trough shown in section. Some of the fluid is diverted in each of the two directions. If the trough is smooth, the velocity of both diverted streams remains $v$, and the only force which can be exerted on the trough is normal to the bottom surface. Hence, the trough will be held in position by forces whose resultant is $F$ normal to the trough. Take $\theta=60^\circ$, $\rho=10^3$ kg/m$^3$, $v=20$ m/s, $A=\sqrt3$ cm$^2$.
\qfig{p2_phy_s3_para5_6_1}
\qn{5} Determine the force $F$ (in N) required to support the trough.
\qn{6} Also find the volume rates of flow $Q_1$ and $Q_2$ for the two streams. What is the ratio $Q_1/Q_2$?
\sect{SECTION-II(ii)}
\qn{7} A thin rod has variable linear mass density given by $\lambda=\lambda_0\left(1+\dfrac{x^2}{L^2}\right)$. Find the distance of its center of mass from the origin (in cm). Take $L=8$ m.
\qfig{p2_phy_s4_q07_1}
\qn{8} A homogeneous circular disk of weight $W=mg$ is suspended from a pin-supported bar (of negligible mass). Initially the disk rotates with the angular velocity $\omega_0$ about its centre. There is a catch in the bar which engages the disc and the disc comes to rest with respect to the bar. What should be $\omega_0$ (in rad/s) so that the system may complete the vertical circle. (Take $\ell=40$ cm, $r=20$ cm, $m=1$ kg, $g=9.8$ m/s$^2$)
\qfig{p2_phy_s4_q08_1}
\qn{9} The wheel shown in Figure rolls with slipping in such a manner that its angular velocity is 5 rad/s in the clockwise direction while the velocity of the centre C is 2 m/s to the left. Determine the speed of point A (in cm/s).
\qfig{p2_phy_s4_q09_1}
\partsec{2}{CHEMISTRY}
\sect{SECTION-I(i)}
\qn{1} A substance which causes pollution is known as pollutant, then which of the following is/are true about pollutant and its effects.
\begin{qoptions}
\opt{A}{SO$_2$ and NO$_2$ are major contributors to acid rain in the form of H$_2$SO$_4$ and HNO$_3$}
\opt{B}{Chemical CF$_2$Cl$_2$ is common component of photochemical smog.}
\opt{C}{Bacteria, Fungi, Moulds , Algae are example of viable particulants.}
\opt{D}{CO$_2$ , CH$_4$ , O$_3$ , water vapour in the atomsphere causes global warming.}
\end{qoptions}
\qn{2} Hydrogen is first element in periodic table , it may exist in the form of hydride, dihydrogen and other forms. Then which of the following statement(s) is/are correct regarding compounds of hydrogen element.
\begin{qoptions}
\opt{A}{Dihydrogen commercially produced by electrolysis of acidified water using platinum electrodes}
\opt{B}{Ionic hydrides are good reducing agents}
\opt{C}{In D$_2$O total number of electrons are greater than H$_2$O}
\opt{D}{Fe$_2$O$_3$ and Cr$_2$O$_3$ are used as catalyst in Bosch process}
\end{qoptions}
\qn{3} (I) K$_3$[Fe(CN)$_6$] \ and (II) K$_4$[Fe(CN)$_6$] reagents are used in qualitative analysis of some of the basic radicals.\par
Then select the \textbf{INCORRECT} option(s).
\begin{qoptions}
\opt{A}{Both the compounds follow Sidgwick E.A.N. rule}
\opt{B}{Compound-I is diamagnetic where as Compound-II is paramagnetic in nature.}
\opt{C}{Both compounds have d$^2$sp$^3$ hybridization with respect of central metal ion.}
\opt{D}{Compound-I is more stable than Compound-II}
\end{qoptions}
\qn{4} (I) [Mn(CO)$_6$]$^+$ , (II) [Cr(CO)$_6$] , (III) [V(CO)$_6$]$^-$\par
are examples of organometallic compounds.\par
Select \textbf{CORRECT} option(s) regarding above given compounds.
\begin{qoptions}
\opt{A}{All are isoelectronic carbonyl compounds}
\opt{B}{In Compound-III carbon-oxygen bond length is maximum}
\opt{C}{[Mn(CO)$_6$] act as reducing agent as compared to [Mn(CO)$_6$]$^+$}
\opt{D}{Metal-Carbon bond has least $\pi$-character in (III) among the Metal-Carbon bond of I, II and III}
\end{qoptions}
\qn{5} To get pure metal from impure metals, there are many refining processes according to nature of metal. Select \textbf{CORRECT} option(s) regarding refining of metals.
\begin{qoptions}
\opt{A}{Poling and electrolytic refining process can be used for refining of impure Cu.}
\opt{B}{Poling , Liquation , Electrolytic refining is connected with refining of Sn metal}
\opt{C}{Zone refining is based upon fractional crystallization, as the impurity prefers to stay in the melt. And it is used for refining of Si , Ge elements.}
\opt{D}{Bassemerization process is used for refining of pig iron to form cast iron}
\end{qoptions}
\qn{6} Concentration of ore is necessary to remove unwanted particles , so which of the following option is/are \textbf{INCORRECT} in terms of concentration methods.
\begin{qoptions}
\opt{A}{In leaching of red bauxite, hall process is used to remove impurity of Fe$_2$O$_3$ and NaOH solution is used as a leaching agent}
\opt{B}{Magnetic separation is used to remove non-magnetic impurity of Wolframite from magnetic casseterite}
\opt{C}{Hydraulic washing is based on the different chemical properties of the ore and gangue particles}
\opt{D}{Potassium or sodium ethyl Xanthate is mixed in frothing mixture as a froth stabilizers in froth floatation method}
\end{qoptions}
\sect{SECTION-I(ii)}
\paraq{7}{8}
Werner was able to explain the nature of bonding in complex , and he concluded that in complex the metal show two different valency\par
(1) Primary valency\par
(2) Secondary valency
\qn{7} For one of the isomer of PtCl$_2$I$_2$.4NH$_3$ , If we add excess AgNO$_3$ solution in 1 mole isomer of PtCl$_2$I$_2$.4NH$_3$ , then formation of silver halide precipitate take place , which is completely insoluble in NH$_3$ solution then formula of isomeric complex is.
\begin{qoptions}
\opt{A}{[Pt(NH$_3$)$_4$I$_2$]Cl$_2$}
\opt{B}{[Pt(NH$_3$)$_4$Cl$_2$]I$_2$}
\opt{C}{[Pt(NH$_3$)$_4$C$\ell$I]C$\ell$I}
\opt{D}{[Pt(NH$_3$)$_5$Cl$_2$]I$_2$}
\end{qoptions}
\qn{8} Which of the following statement is incorrect regarding Co(Br)(Cl)(I).4NH$_3$.
\begin{qoptions}
\opt{A}{3 ionization isomers are possible}
\opt{B}{IUPAC name of the one of the ionization isomers is tetraamminebromidochloridocobalt(III) iodide}
\opt{C}{Primary valency of central metal ion is 3}
\opt{D}{Secondary valency is satisfied by only NH$_3$.}
\end{qoptions}
\paraq{9}{10}
In extraction of gold and silver, the respective metal is leached with dilute solution of NaCN or KCN in the presence of O$_2$. Then the metal is obtained by replacement reaction.
\qn{9} Find out n-factor of Oxygen (O$_2$) and Gold (Au) respectively in reaction in which Gold complex is formed :-
\begin{qoptions}
\opt{A}{4, 1}
\opt{B}{4, 3}
\opt{C}{2, 1}
\opt{D}{1, 3}
\end{qoptions}
\qn{10} Which of the following is \textbf{NOT} related with metallurgy of Au or Ag.
\begin{qoptions}
\opt{A}{Zn act as reducing agent and air act as oxidizing agent}
\opt{B}{Both metal and carbon are sp hybridize in the complex of Ag and Au}
\opt{C}{Complex of Zn, Ag and Au have tetrahedral, linear and linear geometry at metal respectively}
\opt{D}{In the electrolytic refining of Silver, copper metal is obtain in anode mud but Zn metal is not}
\end{qoptions}
\sect{SECTION-II(i)}
\paraq{1}{2}
The crystal field theory (CFT) is an electrostatic model, which considers the metal-ligand bond to be ionic, arising purely from electrostatic interactin between the metal ion and the ligand.
\qn{1} For complex [Co(CN)$_6$]$^{-4}$\par
then find sum of X + Y + Z (by considering C.F.T. only)\par
X = Number of electron in t$_{2g}$ orbital of central metal atom\par
Y = Value of spin only magnetic moment\par
Z = Number of electron in e$_g$ orbital of central metal atom.
\qn{2} Experimental enthalpy of hydration of Cr$^{2+}$ is $-460$ k cal mol$^{-1}$. In the absence of crystal field splitting the value of enthalpy of hydration of Cr$^{2+}$ is $-436$ k cal mol$^{-1}$.\par
(Use 1k cal mol$^{-1}$ = 350 cm$^{-1}$)\par
Find the value of $\left(\dfrac{\Delta}{10}\right)$ for [Cr(H$_2$O)$_6$]$^{2+}$ in cm$^{-1}$ :
\paraq{3}{4}
\qfig{p2_chem_s3_para3_4_1}
The free energy changes that occur when one gram molecule of a common reactant (in this case dioxygen) is used may be plotted graphically against temperature for a number of reactions of metals of their oxides. This graph is shown in figure and is called an Ellingham diagram (for oxides).
\qn{3} Select the number of \textbf{CORRECT} statement(s) about the plot is / are:
\qfig{p2_chem_s3_q03_1}
\centerline{\small Temperature}
\par\noindent(i) T$_1$ and T$_2$ are melting point \& boiling point of Mg respectively.\par
(ii) T$_1$ and T$_2$ are melting point \& boiling point of MgO respectively.\par
(iii) Reduction of MgO by coke is possible above T$_3$\par
(iv) Mg can be extracted from gaseous products by rapid cooling.\par
(v) After T$_1$ point $|\Delta G^\circ|$ decreases\par
(vi) After T$_2$ point $|\Delta G^\circ|$ increases
\qn{4} As per given ellingham diagram (in paragraph) at 1500 K, if $\Delta G^\circ$ for the oxidation of Si is $-700$ kJ/mole of O$_2$ and $\Delta G^\circ$ for the reduction of MgO is $+850$ kJ/mole of O$_2$.\par
Find $|\Delta G^\circ|_{\text{redox}}$ for the reaction of\par
2MgO + Si $\rightleftharpoons$ SiO$_2$ + 2Mg at 1500 K
\paraq{5}{6}
Compounds which have the same molecular formula, but differ in their properties due to the difference in structure or orientation in 3D, are known as isomers.
\qn{5} A complex [CoBr(ONO)(en)$_2$]Cl.H$_2$O, can show how many types of isomerism from given.\par
Linkage , Coordination , Ionization , Geometrical , Hydrate , Optical , Stereo isomerism
\qn{6} The octahedral complex [Co(en)$_2$(NO$_2$)$_2$]$^+$ can exist in total number of isomeric form including structural \& stereo isomers is/are :-
\sect{SECTION-II(ii)}
\qn{7} How many statements is/are \textbf{CORRECT} about H-atom and its compound(s).\par
(i) Ortho H$_2$ $>$ Para H$_2$ : Stability at low temperature\par
(ii) H$_2$ $<$ D$_2$ $<$ T$_2$ : Boiling point order\par
(iii) Atomic hydrogen $>$ Nascent hydrogen $>$ Molecular hydrogen : Reactivity order\par
(iv) Diborane is electron$-$precise covalent hydride\par
(v) Temporary hardness is due to presence of soluble salts of magnesium and calcium in the form of chlorides and sulphates\par
(vi) Equimolar mixture of CO and H$_2$ is known as water gas.
\qn{8} Pollution of water originates from Human activities. Select \textbf{CORRECT} points related to water pollution.\par
(1) Clean water would have BOD value of less than 5 ppm where as highly pollutated water would have a BOD value of 17 ppm or more.\par
(2) F$^-$ concentration above 2 ppm cause brown mottling of teeth.\par
(3) Upper limit concentration of lead in drinking water is about 50 ppb\par
(4) Maximum limit of nitrate in drinking water is 50 ppm.\par
(5) Greater than 500 ppm of sulphate in drinking water causes Laxative effect\par
(6) 0.2 ppm of Fe and Al and 3 ppm of Cu metal is maximum prescribed concentration in drinking water.
\qn{9} In the following ores find sum of (x + y + z) where\par
x = number of ores in which two metal atoms exist in different oxidation state.\par
y = number of ores in which two different metals are available.\par
z = number of ores in which atleast one metal ion is diamagnetic in nature.\par
Magnetite, Chromite, Limonite, Malachite, I$\ell$menite, Chalcopyrite, Dolomite, Azurite, Carnalite, Cryolite
\partsec{3}{MATHEMATICS}
\sect{SECTION-I(i)}
\qn{1} If \[A=\begin{bmatrix}\scriptstyle x+xy+xyz&0&0\\0&\scriptstyle y+yz+xyz&0\\0&0&\scriptstyle z+zx+xyz\end{bmatrix}\] $|A|\neq0$ and $xyz=1$ and sum of maximum and minimum value of $\text{Tr}(A^{-1})$ (where $\text{Tr}(A)$ denotes sum of diagonal elements) is $\dfrac pq$, then which of following options is/are correct ($p$ and $q$ are coprime)
\begin{qoptions}
\opt{A}{$p+q=3$}
\opt{B}{$p+q=5$}
\opt{C}{$p-q=1$}
\opt{D}{$3p-2q=4$}
\end{qoptions}
\qn{2} Let there are 6 hats of different colours and 6 shirts of same colours as that of hats. If the number of ways in which these can be put on by 5 men such that no man wear the shirt and the hat of the same colour is $K\times6!$. Then the value of $K$ is greater than
\begin{qoptions}
\opt{A}{108}
\opt{B}{218}
\opt{C}{258}
\opt{D}{309}
\end{qoptions}
\qn{3} Let $f(x)$ is a polynomial in $x$ defined as $f(x)=(1+x+x^2)(1+x)^5(1+x^2)^5$ then -
\begin{qoptions}
\opt{A}{coefficient of $x^{17}$ in $f(x)$ is 1}
\opt{B}{coefficient of $x^{16}$ in $f(x)$ is 6}
\opt{C}{coefficient of $x^{15}$ in $f(x)$ is 11}
\opt{D}{coefficient of $x$ in $f(x)$ is 6}
\end{qoptions}
\qn{4} Let $B=[b_{ij}]_{n\times n}$ be a square matrix of order $n$ such that $b_{ij}=\begin{cases}0&;\ i\neq j\\ \dfrac{4+i+j}{2}&;\ i=j\end{cases}$
\begin{qoptions}
\opt{A}{$\displaystyle\lim_{m\to\infty}\left(\dfrac{\sum_{n=1}^{m}\text{trace}(B)}{m^3}\right)=\dfrac16$}
\opt{B}{$\displaystyle\sum_{n=1}^{10}\text{trace}(B)=330$}
\opt{C}{$\displaystyle\sum_{n=1}^{10}\text{trace}(B)=320$}
\opt{D}{$\displaystyle\lim_{m\to\infty}\left(\dfrac{\sum_{n=1}^{m}\text{trace}(B)}{m^3}\right)=\dfrac13$}
\end{qoptions}
\qn{5} $D=\begin{vmatrix}a&1&1\\1&1&b\\c&1&1\end{vmatrix}\times\begin{vmatrix}1-b&bc-1&1-c\\0&a-c&c-a\\b-1&1-ab&a-1\end{vmatrix}$. If $(a-c)(1-b)=5$, then value of $D$ will be -
\begin{qoptions}
\opt{A}{An even number}
\opt{B}{An odd number}
\opt{C}{A prime number}
\opt{D}{A composite number}
\end{qoptions}
\qn{6} Let $P$ is a skew symmetric matrix of order $3\times3$ with real entries. Given $I+P$ and $I-P$ ($I$ is identity matrix of order $3\times3$) are non-singular matrices and $Q=(I+P)(I-P)^{-1}$ (Where $\det(Q)>0$).\par
Identify the correct options.
\begin{qoptions}
\opt{A}{$\det(2Q)=8$}
\opt{B}{$QQ^T=I$}
\opt{C}{$\det(\text{adj}Q)=1$}
\opt{D}{$Q=Q^{-1}$}
\end{qoptions}
\sect{SECTION-I(ii)}
\paraq{7}{8}
$g(y)=\begin{vmatrix}p-y&q&r\\r&p-y&q\\q&r&p-y\end{vmatrix}$, where $p,q,r$ are distinct positive numbers\par
On the basis of above information answer the following :
\qn{7} Absolute term in $g(y)$ can be -
\begin{qoptions}
\opt{A}{2}
\opt{B}{$-1$}
\opt{C}{0}
\opt{D}{All of the above}
\end{qoptions}
\qn{8} Coefficient of $y$ in $g(y)$ will be -
\begin{qoptions}
\opt{A}{independent of only $p$}
\opt{B}{independent of $q$ as well as $r$}
\opt{C}{independent of only $r$}
\opt{D}{none of these}
\end{qoptions}
\paraq{9}{10}
Let $A=\begin{bmatrix}1&2&3\\0&2&3\\0&0&3\end{bmatrix}$ and $U_1$, $U_2$ and $U_3$ are first, second and third columns of a $3\times3$ matrix $U$ respectively.\par
If column matrices $U_1$, $U_2$ and $U_3$ satisfy $AU_1=\begin{bmatrix}1\\0\\0\end{bmatrix}$, $AU_2=\begin{bmatrix}0\\1\\0\end{bmatrix}$, $AU_3=\begin{bmatrix}0\\0\\1\end{bmatrix}$.\par
On the basis of above information, answer the following questions :
\qn{9} The value $|U|$ is
\begin{qoptions}
\opt{A}{$\dfrac16$}
\opt{B}{$-\dfrac13$}
\opt{C}{$\dfrac13$}
\opt{D}{$\dfrac12$}
\end{qoptions}
\qn{10} $U^{-1}$ is a
\begin{qoptions}
\opt{A}{Symmetric matrix}
\opt{B}{Skew symmetric matrix}
\opt{C}{Triangular matrix}
\opt{D}{Scalar matrix}
\end{qoptions}
\sect{SECTION-II(i)}
\paraq{1}{2}
Let $A6P+BQ4+C45=887$, where `XYZ' denotes three digits number.\par
On the basis of above information, answer the following questions :
\qn{1} If the number of ordered triplets $(A,B,C)$ is $\lambda$ then value of $\lambda/2$ is
\qn{2} If the number of solutions that product of $A$, $B$ and $C$ is an odd integer is $\mu$ then value of $\mu/5$ is
\paraq{3}{4}
$n$ straight lines are drawn in a plane such that no 2 of them are parallel and no 3 of them are concurrent.\par
Let $T_n$ represents the number of parts in which the $n$ lines divides the plane and let $Q_n=\displaystyle\prod_{r=1}^{n}\dfrac{T_{r+1}-1}{T_r-1}$.
\qn{3} $T_{10}$ is equal to
\qn{4} $\displaystyle\sum_{r=1}^{10}\dfrac{6}{Q_r}$ is equal to
\paraq{5}{6}
If $a_n=\displaystyle\sum_{r=0}^{n}\dfrac1{{}^nC_r}$ then $\displaystyle\sum_{r=0}^{n}\dfrac{r^2}{{}^nC_r}=P(n)a_{n+2}+Q(n)a_{n+1}+a_n+R(n)$\par
Where $P(n)$, $Q(n)$, $R(n)$ are polynomial functions of $n$ then
\qn{5} $\dfrac{15}{10-R(4)}$ is equal to
\qn{6} $P(5)+Q(5)$ is equal to
\sect{SECTION-II(ii)}
\qn{7} If $A$ \& $B$ are two square matrices of order $3\times3$ such that $2A+3BB^T=I$ and $B^{-1}=A^T$.\par
Let $\alpha=\det(A^{-1}-3B^3+BA)$ and $\beta=\text{Tr}(A^{-1}-3B^3-BA)$ then value of $\left(\dfrac\alpha\beta\right)$ is
\qn{8} The value of $\displaystyle\lim_{n\to\infty}\sum_{r=1}^{n}\left(\sum_{t=0}^{r-1}\dfrac1{5^n}\,{}^nC_r\,{}^rC_t\,3^t\right)$ is equal to
\qn{9} If $a$, $b$, $c$ be the distinct real numbers and $x$, $y$, $z$ be any three variables such that the system of linear equations\par
$(3z-2y)a+(2x-z)b+(y-3x)c=0$\par
$(y-3x)a+(3z-2y)b+(2x-z)c=0$\par
$(2x-z)a+(y-3x)b+(3z-2y)c=0$\par
can be express as $AX=O$, $A$ is $3\times3$ square matrix and $X=\begin{bmatrix}x\\y\\z\end{bmatrix}$ then the value of $\det(A)+\det(2A)+\det(3A)$ is :

\end{multicols}
\end{document}
